\documentclass[12pt]{article}
\usepackage[margin=1in]{geometry}
\usepackage{tabularx}
\usepackage{booktabs}
\usepackage{hyperref}
\usepackage{parskip}
\usepackage{enumitem}

\hypersetup{
    colorlinks=true,
    linkcolor=blue,
    urlcolor=blue
}

\title{\textbf{Design Specifications Document} \\ EagleNav -- Campus Navigation Using Augmented Reality}
\author{
Jian Mejia, Lcndr Aquino, George Barrone, Abigail Diaz, \\
Isidoro Flores, Santhosh Kumar, David Lopez, Reda Masri, Athena Ruiz \\
\\
California State University, Los Angeles \\
Office for Students with Disabilities (OSD)
}
\date{May 2026}

\begin{document}

\maketitle

\begin{center}
\begin{tabularx}{0.8\textwidth}{|l|X|}
\hline
\textbf{Version} & 1.0 \\
\hline
\textbf{Date} & May 2026 \\
\hline
\textbf{Sponsor} & Office for Students with Disabilities (OSD) \\
\hline
\textbf{Faculty Advisor} & Dr. David Krum \\
\hline
\textbf{Platform} & Android 10+ / iOS 14+ | Flutter (Dart) \\
\hline
\textbf{Status} & Final -- Academic Year 2025--2026 \\
\hline
\end{tabularx}
\end{center}

\newpage
\tableofcontents
\newpage

\section{Project Overview}

\subsection{Purpose}
EagleNav is an augmented reality campus navigation platform built for Android and iOS that uses custom high-precision mapping of the Cal State LA campus. The application is designed to support students with disabilities by offering accessible, independence-focused navigation tools including AR directional arrows, voice-guided directions, LiDAR-enhanced obstacle awareness, and seamless integration with campus events and building data. EagleNav aims to remove navigation barriers and create a more inclusive campus experience for all students, faculty, staff, and visitors at California State University, Los Angeles.

\subsection{Scope}
\begin{itemize}
    \item Geographic coverage: Cal State LA main campus only -- academic buildings, student services, housing perimeter paths, and all parking structures.
    \item Indoor coverage (v1): Major public corridors, elevators, ramps, and primary entrances/exits of high-traffic buildings.
    \item Off-campus routing is not supported.
\end{itemize}

\subsection{Project Team}
\begin{tabularx}{\textwidth}{|l|l|X|}
\hline
\textbf{Name} & \textbf{Sub-team} & \textbf{Role} \\
\hline
Jian Mejia & Computer Vision & Computer Vision Lead \\
\hline
Isidoro Flores & Computer Vision & Computer Vision \\
\hline
Santhosh Kumar & Computer Vision & Computer Vision \\
\hline
Abigail Diaz & Path Tracing & Path Tracing Lead \\
\hline
Reda Masri & Path Tracing & Path Tracing \\
\hline
George Barrone & Path Tracing & Path Tracing \\
\hline
Lcndr Aquino & Haptics & Haptics \\
\hline
Athena Ruiz & Haptics & Haptics \\
\hline
David Roberto Lopez & Data & Data \\
\hline
\end{tabularx}

\subsection{Key Stakeholders}
\begin{itemize}
    \item Project Sponsor: Office for Students with Disabilities (OSD), Cal State LA
    \item Faculty Advisor: Dr. David Krum
    \item Primary Users: CSULA students with disabilities (wheelchair, visually impaired, hearing impaired)
    \item Secondary Users: General student body, OSD staff, campus security, campus visitors
\end{itemize}

\section{Design Goals and Principles}

\subsection{Core Design Goals}
\begin{itemize}
    \item \textbf{Accessibility First:} All features must be usable by students with visual, hearing, mobility, and cognitive disabilities.
    \item \textbf{Independence-Focused:} Navigation should empower users to navigate without requiring assistance from others.
    \item \textbf{Multimodal Feedback:} Every critical interaction must provide at minimum two of: visual, audio, and haptic feedback.
    \item \textbf{Graceful Degradation:} The app must function in degraded connectivity environments and fall back gracefully.
    \item \textbf{Minimal Cognitive Load:} UI complexity is strictly limited; interfaces must be intuitive without training.
\end{itemize}

\subsection{Design Principles}
\begin{itemize}
    \item Deterministic routing: identical inputs always produce identical routes.
    \item Body-relative directions: all navigation cues reference the user's current facing direction, not cardinal directions.
    \item Redundant signals: every turn or alert is communicated via color + vibration + text.
    \item Independent voice channel: app speech is fully separate from system accessibility tools (VoiceOver / TalkBack).
    \item No gesture conflicts with OS accessibility layers.
\end{itemize}

\section{System Architecture}

\subsection{Architectural Overview}
EagleNav follows a three-tier layered architecture: Cloud Infrastructure, an on-device Controller Layer, and an on-device Service Layer. Each tier has a single responsibility, communicates only downward through well-defined interfaces, and is independently replaceable.

\subsection{Cloud Infrastructure}
\begin{itemize}
    \item \textbf{Cloudflare R2:} Object storage hosting the campus OpenStreetMap extract (\texttt{renum\_file.osm.pbf}). Valhalla loads this file into memory on startup.
    \item \textbf{Valhalla Routing Engine:} Open-source routing engine running on a Google Cloud Platform virtual machine; exposes a REST API consumed by the Flutter client.
\end{itemize}

\subsection{Controller Layer (On-Device)}
\begin{tabularx}{\textwidth}{|l|X|}
\hline
\textbf{Controller} & \textbf{Responsibility} \\
\hline
LocationController & Manages GPS stream, filters fixes by accuracy, exposes reactive current position. \\
\hline
RoutingController & Fetches routes from Valhalla, tracks deviation, triggers reroutes. \\
\hline
GuidanceController & Tracks step progression, advances steps at turn points, detects arrival. \\
\hline
NavigationVoiceController & Produces spoken instructions, manages TTS priority, exposes smoothed compass heading. \\
\hline
UiNavigationController & Owns four UI states: idle, destinationSelected, navigating, arrived. \\
\hline
\end{tabularx}

\subsection{Service Layer (On-Device)}
\begin{tabularx}{\textwidth}{|l|X|}
\hline
\textbf{Service} & \textbf{Responsibility} \\
\hline
LocationService & Wraps geolocator plugin; exposes filtered position stream. \\
\hline
ValhallaService & Formats routing requests, posts over HTTPS, parses JSON response into typed step objects. \\
\hline
CompassService & Wraps flutter\_compass; smooths magnetometer signal; exposes stable heading. \\
\hline
TTS Manager & Wraps flutter\_tts with priority queue -- urgent cues interrupt non-urgent ones. \\
\hline
HapticFeedback & Wraps device vibration API; produces patterned cues for turn events. \\
\hline
InstructionEnhancer & Converts Valhalla maneuver codes, distances, and angles into body-relative spoken phrases. \\
\hline
OrientationCoach & At route start, compares compass heading to first-segment bearing; produces alignment cues. \\
\hline
OffCourseDetector & Uses heading-aware thresholds (45m when aligned, 15m when opposed) for reroute decisions. \\
\hline
LandmarkService & On user request, computes nearby buildings and speaks their body-relative position. \\
\hline
\end{tabularx}

\subsection{Technology Stack}
\begin{tabularx}{\textwidth}{|l|X|}
\hline
\textbf{Layer / Component} & \textbf{Technology} \\
\hline
Mobile Framework & Flutter 3.x / Dart 3.x (cross-platform iOS \& Android) \\
\hline
Maps & flutter\_map with OpenStreetMap tiles \\
\hline
Routing Engine & Valhalla (open-source) on Google Cloud Platform via REST \\
\hline
Map Data & OpenStreetMap extract (.osm.pbf) on Cloudflare R2 \\
\hline
AR & ARCore (Android) / ARKit (iOS) via Flutter platform channels \\
\hline
Computer Vision (Android) & YOLO-based object detection via TensorFlow Lite \\
\hline
Computer Vision (iOS) & YOLO-based with CoreML; LiDAR hardware depth on supported devices \\
\hline
Local Storage & SharedPreferences (settings/bookmarks); JSON/GeoJSON assets \\
\hline
Security & TLS for all network calls; PII minimized; location processed on-device when possible \\
\hline
\end{tabularx}

\section{Data Flow Design}

\subsection{Route Planning Subsystem}
\begin{tabularx}{\textwidth}{|l|l|X|}
\hline
\textbf{\#} & \textbf{Process} & \textbf{Description} \\
\hline
P1.1 & Building Selector & Resolves closest accessible entrance from buildings.json using current GPS. \\
\hline
P1.2 & RoutingController & Assembles routing request and manages request lifecycle including reroutes. \\
\hline
P1.3 & ValhallaService & Formats request as JSON, POSTs to Valhalla over HTTPS, parses response into typed route object. \\
\hline
P1.4 & Valhalla Engine & Computes optimal pedestrian route with sidewalk preference and unpaved exclusion. \\
\hline
\end{tabularx}

\subsection{Guidance Subsystem}
LocationController emits a filtered GPS position which the screen forwards to RoutingController and GuidanceController. RoutingController computes deviation using heading-aware thresholds (45m when aligned, 15m when opposed). GuidanceController checks step progression and advances the step index when the user passes a step endpoint. InstructionEnhancer converts maneuver codes and compass angles into body-relative spoken phrases, and HapticFeedback fires alongside each spoken cue.

\subsection{Compass and Orientation Subsystem}
The CompassService smooths magnetometer readings and exposes \texttt{currentHeading} as a reactive notifier. The OrientationCoach runs at navigation start and selects a coaching phrase based on the angle between the user's heading and the route's first segment:

\begin{tabularx}{\textwidth}{|l|X|}
\hline
\textbf{Angle Range} & \textbf{Coaching Phrase} \\
\hline
$|$angle$|$ $<$ 15 degrees & Walk straight ahead \\
\hline
15 -- 45 degrees & Slightly to your right / left \\
\hline
45 -- 150 degrees & Turn to your right / left \\
\hline
$|$angle$|$ $>$ 150 degrees & Turn around \\
\hline
\end{tabularx}

\subsection{Landmark Service Subsystem}
The Landmark Service is user-initiated and operates independently of active navigation. It finds nearby buildings within a 50m radius, computes a body-relative bearing for each, converts metric distance to approximate step count (0.76m stride), and announces the result (e.g., ``Library is to your front right, about 30 steps away.'').

\subsection{Events Subsystem}
The EventsScreen fetches HTML from the CSULA Events website, parses it into structured event data, and displays it to the user. Users can bookmark events stored locally via SharedPreferences, and a notification scheduler fires local notifications for upcoming bookmarked events.

\subsection{Computer Vision Subsystem}
The camera captures real-time video input processed by YOLO-based detection and segmentation models. Detected objects are filtered by direction (left, ahead, right) and urgency, then delivered via TTS and haptics. Android uses a TensorFlow Lite depth estimator; iOS uses LiDAR hardware on supported devices with CoreML model equivalents otherwise.

\section{Functional Requirements}

\subsection{General Application Behavior}
\begin{itemize}
    \item FR-0.1: Display the campus map as the default view on launch.
    \item FR-0.2: Provide access to Home, Favorites, Alerts, Profile, and Emergency Assistance from all views.
    \item FR-0.3: On launch, focus the search bar and announce via TTS: ``Search for a destination.''
\end{itemize}

\subsection{Destination Selection and Confirmation}
\begin{itemize}
    \item FR-1.1 / 1.2: Search for buildings by typing or speaking; results update dynamically.
    \item FR-2.1: Display generated route color-coded before navigation begins.
    \item FR-2.3: Display route summary: total distance, estimated travel time, and entrance name.
    \item FR-2.5: Show ``Start Navigation'' button after route summary is presented.
\end{itemize}

\subsection{Core Navigation}
\begin{itemize}
    \item FR-3.1: Use a custom pedestrian path network for on-campus routing (Valhalla with campus OSM data).
    \item FR-3.2: At navigation start, provide turn-by-turn guidance and announce initial walking direction.
    \item FR-3.3: Alert the user of obstacles (stairs, blockages, construction).
    \item FR-3.4: Auto-recalculate route when the user deviates from the planned path.
    \item FR-3.5: Vibrate when off-route or approaching a turn.
\end{itemize}

\subsection{Navigation Enhancements (v1.6)}
\begin{itemize}
    \item FR-14.1: Route to the closest accessible entrance from buildings.json.
    \item FR-14.2: Express all directions body-relative (e.g., ``slightly to your right''), not cardinal.
    \item FR-14.3: Apply heading-aware deviation thresholds (45m when aligned, 15m when opposed).
    \item FR-14.4: Two-phase navigation start -- Phase 1 guides to path network; Phase 2 begins turn-by-turn.
    \item FR-14.5: TTS priority queue -- urgent cues (turn-now, arrival) interrupt non-urgent ones.
    \item FR-14.7: Express distances in approximate step counts, not meters.
    \item FR-14.8: Arrival detection at 8m from final step endpoint; announce door direction body-relative.
    \item FR-14.9: Valhalla pedestrian costing: prefer sidewalks/walkways, penalize alleys/tunnels, exclude unpaved.
\end{itemize}

\subsection{Accessibility Requirements}

\subsubsection{Mobility Support}
Large, high-contrast touch targets throughout. Voice commands for destination entry and navigation start. One-tap shortcuts to library, cafeteria, and OSD office. Prioritize ramps, elevators, and automatic doors in route generation.

\subsubsection{Hearing Support}
Distinct haptic patterns for left turn, right turn, straight, and wrong turn. Text mode for turn-by-turn instructions as readable text. Customizable vibration intensity and distance alerts.

\subsubsection{Vision Support}
Full compatibility with VoiceOver (iOS) and TalkBack (Android). Spoken descriptions of all menus, buttons, and navigation steps. Adjustable font sizes and zoom levels.

\subsubsection{Cognitive Support}
Simple, sequential navigation instructions. Landmark-based guidance (e.g., ``Turn right at the fountain''). Optional confirmation after each instruction.

\section{Non-Functional Requirements}

\begin{tabularx}{\textwidth}{|l|l|X|}
\hline
\textbf{Category} & \textbf{Requirement} & \textbf{Target} \\
\hline
Performance & Route generation latency & 95\% of routes computed in less than 3 seconds \\
\hline
Performance & App cold start & Less than 2.5 seconds \\
\hline
Performance & UI frame budget & 16ms target (60 fps) \\
\hline
Reliability & Core function uptime & 99\% uptime \\
\hline
Safety & Emergency button offline & Must function offline via SMS fallback \\
\hline
Security & Data encryption & End-to-end TLS; location data encrypted at rest and in transit \\
\hline
Security & PII minimization & No student records; location processed on-device when possible \\
\hline
Usability & Accessibility default & Accessibility features enabled on first run \\
\hline
Compatibility & OS support & Android 10+ and iOS 14+ \\
\hline
\end{tabularx}

\section{User Interface Design}

\subsection{Design Principles}
\begin{itemize}
    \item Clean, minimal UI -- reduces cognitive load and makes core functions immediately discoverable.
    \item Large touch targets -- minimum 44x44pt (Apple HIG) for low-vision users.
    \item Optional high-contrast mode toggled from the accessibility profile.
    \item Customizable font sizes from the settings screen.
    \item Screen reader support as a first-class requirement, not an afterthought.
\end{itemize}

\subsection{Screen Inventory}
\begin{tabularx}{\textwidth}{|l|X|}
\hline
\textbf{Screen} & \textbf{Key Elements} \\
\hline
Home / Map & Search bar, campus map, current-location pin, bottom navigation bar \\
\hline
Route Preview & Color-coded route polyline, destination marker, route summary card, Start Navigation button \\
\hline
Navigation / Guidance & Turn banner, progress indicator, AR button, haptic/voice toggles, End/Cancel button \\
\hline
Favorites & Saved locations, quick-start navigation, edit/delete options \\
\hline
Events & Campus event list, bookmark toggle, local notification scheduler \\
\hline
Accessibility Profile & Profile selector, TTS speed/pitch/volume, haptic intensity, high-contrast toggle, font size slider \\
\hline
Emergency & One-touch campus security call, Safe Walk toggle, confirmation prompt \\
\hline
\end{tabularx}

\subsection{UI Navigation Flow}
HomeMap $\rightarrow$ (Search or Favorite tap) $\rightarrow$ RoutePreview $\rightarrow$ Start $\rightarrow$ Guidance $\rightarrow$ (Deviation? $\rightarrow$ Recalculate) $\rightarrow$ Arrive $\rightarrow$ Done

\section{Database and Data Design}

\subsection{Local Data Assets}
\begin{tabularx}{\textwidth}{|l|X|}
\hline
\textbf{File / Store} & \textbf{Purpose} \\
\hline
buildings.json & Campus building names with accessible entrance coordinates. \\
\hline
buildings.geojson & Building footprints and entrance coordinates for map rendering. \\
\hline
bookmarks.json & User-saved routes and locations. \\
\hline
SharedPreferences & User accessibility preferences, app settings, and bookmarked events. \\
\hline
depth\_estimator.tflite & Android TensorFlow Lite model for object distance estimation. \\
\hline
object\_detector.tflite & Android TensorFlow Lite model for real-time YOLO object detection. \\
\hline
labels.txt & Classification label definitions for the object detection model. \\
\hline
renum\_file.osm.pbf & Campus OpenStreetMap extract on Cloudflare R2. \\
\hline
\end{tabularx}

\subsection{Data Security Constraints}
\begin{itemize}
    \item All network calls over TLS; no plaintext transmission of location data.
    \item No student records stored or transmitted; PII minimized throughout.
    \item Sensitive data encrypted at rest on-device.
    \item Emergency location sharing transmits only to campus security / OSD, not third parties.
\end{itemize}

\section{Implementation Status (v1.6)}

\begin{tabularx}{\textwidth}{|l|l|l|X|}
\hline
\textbf{FR} & \textbf{Feature Group} & \textbf{Status} & \textbf{Notes} \\
\hline
FR-0 & General Application Behavior & Implemented & Map default view, core navigation, TTS on launch. \\
\hline
FR-1 & Destination Selection & Implemented & Search, dynamic results, map tap, clear. \\
\hline
FR-2 & Destination Confirmation & Implemented & Route preview, summary, Start Navigation. \\
\hline
FR-3 & Core Navigation & Implemented & Turn-by-turn, deviation detection, rerouting, haptics. \\
\hline
FR-4 & Accessibility -- Mobility & Partial & Voice commands, shortcuts, emergency profiles deferred. \\
\hline
FR-5 & Accessibility -- Hearing & Partial & Haptic turn patterns implemented. Custom settings deferred. \\
\hline
FR-6 & Accessibility -- Vision & Partial & Screen reader, font/zoom implemented. Voice input deferred. \\
\hline
FR-13 & Speech Feedback (TTS) & Implemented & Full TTS pipeline implemented. \\
\hline
FR-14 & Navigation Enhancements & Implemented & Closest-entrance, body-relative, heading-aware, two-phase start. \\
\hline
FR-15 & Orientation \& Awareness & Implemented & CompassService, OrientationCoach, LandmarkService. \\
\hline
FR-9 & Personalization \& Scheduling & Future Work & All requirements deferred. \\
\hline
FR-12 & Accessibility -- Color Blind & Future Work & All color-blind themes deferred. \\
\hline
\end{tabularx}

\section{Testing and Validation Strategy}

\begin{itemize}
    \item \textbf{Unit Tests:} CompassService smoothing, InstructionEnhancer phrase generation, OffCourseDetector thresholds.
    \item \textbf{Integration Tests:} Route planning pipeline; Events parsing and notification scheduling.
    \item \textbf{Walk Tests:} End-to-end turn-by-turn guidance; reroute on deviation; arrival detection; LandmarkService with real buildings.
    \item \textbf{Manual Tests:} TTS priority interruption; VoiceOver / TalkBack compatibility; gesture non-conflict.
    \item \textbf{Performance Tests:} Route latency (less than 3s 95th percentile); cold start (less than 2.5s); UI frame time.
\end{itemize}

\section{Future Work}

\subsection{Deferred Features}
\begin{itemize}
    \item Voice commands for destination entry and navigation initiation.
    \item Elevator/ramp preference in route costing.
    \item Emergency button with offline SMS fallback.
    \item Full accessibility user profiles -- wheelchair, cane, temporary injury.
    \item Color-blind UI themes (protanopia, deuteranopia, tritanopia).
    \item Indoor room-to-room routing and wing/floor search.
    \item Personalization and scheduling -- bookmarks, class schedules, SIS integration.
\end{itemize}

\subsection{Potential Enhancements}
\begin{itemize}
    \item SIS integration to auto-import class schedules and provide next-class routing suggestions.
    \item Multi-language support (en-US baseline with right-to-left readiness).
    \item Wearable device integration for haptic-only navigation guidance.
\end{itemize}

\section{Glossary}

\begin{tabularx}{\textwidth}{|l|X|}
\hline
\textbf{Term} & \textbf{Definition} \\
\hline
AR & Augmented Reality -- overlays digital navigation markers on the physical world via the device camera. \\
\hline
Body-relative direction & Direction expressed relative to the user's current facing rather than cardinal direction. \\
\hline
CompassService & On-device service that smooths and validates raw magnetometer readings, exposing currentHeading. \\
\hline
GuidanceController & Real-time navigation processor tracking step progression and detecting arrival. \\
\hline
InstructionEnhancer & Converts Valhalla maneuver codes, distances, and angles into spoken body-relative phrases. \\
\hline
LandmarkService & Identifies nearby campus buildings and announces body-relative direction and step distance. \\
\hline
OffCourseDetector & Uses heading-aware thresholds to decide when to trigger a route recalculation. \\
\hline
OrientationCoach & One-shot service at navigation start aligning the user with the route's first segment. \\
\hline
OSD & Office for Students with Disabilities at Cal State LA; primary project sponsor. \\
\hline
RoutingController & On-device controller managing route requests, active route state, and rerouting. \\
\hline
TTS & Text-to-Speech -- converts navigation instructions to spoken audio. \\
\hline
Valhalla & Open-source routing engine hosted on GCP, computing pedestrian routes from OSM data. \\
\hline
\end{tabularx}

\section{References}

\begin{enumerate}
    \item EagleNav Software Requirements Specification (SRS) v1.6 -- OSD, Cal State LA, May 2026
    \item EagleNav Software Design Document (SDD) v4.0 -- OSD, Cal State LA, May 2026
    \item IEEE Std 830-1998 -- IEEE Recommended Practice for Software Requirements Specifications
    \item CSULA Campus Map. \url{https://www.calstatela.edu/map}
    \item Valhalla Routing Engine. \url{https://github.com/valhalla/valhalla}
    \item Flutter / Dart Documentation. \url{https://docs.flutter.dev}
    \item Apple Human Interface Guidelines -- Accessibility
    \item Android Accessibility Developer Guide -- TalkBack \& Touch Targets
\end{enumerate}

\end{document}
